\documentclass[11pt]{article}
\usepackage[margin=1in]{geometry}
\usepackage{amsmath,amssymb,amsthm,mathtools}
\usepackage{booktabs}
\usepackage{xcolor}
\usepackage{microtype}
\usepackage[colorlinks=true,linkcolor=blue!55!black,citecolor=blue!55!black,urlcolor=blue!55!black]{hyperref}

\newtheorem{theorem}{Theorem}
\newtheorem{proposition}[theorem]{Proposition}
\newtheorem{lemma}[theorem]{Lemma}
\newtheorem{corollary}[theorem]{Corollary}
\newcommand{\eps}{\{-1,1\}}
\newcommand{\E}{\mathbb E}
\newcommand{\cE}{\mathcal E}
\newcommand{\Ftwo}{\mathbb F_2}
\DeclareMathOperator{\diag}{diag}
\DeclareMathOperator{\tr}{tr}
\DeclareMathOperator{\wt}{wt}

\hypersetup{
  pdftitle={A Second Attempt at MathOverflow 413935},
  pdfauthor={OpenAI Codex},
  pdfsubject={Relaxation theorems and obstructions for a min-max quadratic sign problem},
  pdfkeywords={quadratic forms, signs, discrepancy, spin glasses, cut codes, graphons}
}

\title{\textbf{A Second Attempt at MathOverflow 413935}\\[4pt]
\large Exact relaxations, finite-temperature barriers, and a sharper proof frontier}
\author{OpenAI Codex\\
\small AI-generated research note; directed and preserved by Rob Sneiderman}
\date{1 August 2026}

\begin{document}
\maketitle

\begin{center}
\fcolorbox{red!65!black}{red!4}{
\begin{minipage}{0.89\textwidth}
\textbf{Status and nonclaim.}
This note does not prove that the limit in Question~\ref{q:main} exists, and
it does not prove nonexistence.  It contains independently checked partial
theorems and exact finite computations.  It must not be cited or submitted as
a solution to the MathOverflow question.
\end{minipage}}
\end{center}

\begin{abstract}
For a symmetric zero-diagonal sign matrix \(A\), put
\[
 M(A)=\max_{x\in\eps^n}\left|\frac12x^{\mathsf T}Ax\right|,
 \qquad F(n)=\min_A M(A).
\]
The open question is whether \(F(n)/n^{3/2}\) converges.  We first audit the
previous bounds \(1/\pi\le\liminf\le\limsup\le1/2\), the augmented cut-code
identity, and one-vertex monotonicity; all survive audit.  We then show that
the proposed cube relaxation has the unique optimizer zero at zero temperature
and at every fixed unscaled positive temperature, with a leading-order
zero-temperature integrality gap.  In contrast, an elliptope relaxation has
the exact thermodynamic limit \(1/2\), while the defined two-covariance,
edgewise-oriented Gaussian certificate has asymptotic lower-bound constant
\(1/\pi\).
We identify the augmented cut code as a binary linear code whose dual consists
of even-cardinality Eulerian subgraphs, and derive an exact signed-Eulerian
MacWilliams formula.  Ordinary graphons and limiting empirical spectra are
proved to erase the \(n^{3/2}\) functional.  Exhaustive computation gives
\(F(2),\ldots,F(10)=1,3,4,4,5,9,10,12,13\).  The sharp remaining target is a
quenched minimax free-energy limit that absorbs, rather than separately
bounds, the cross-block contribution.
\end{abstract}

\section{Problem, audited baseline, and scope}
\label{q:main}

The question, asked by Paata Ivanisvili on MathOverflow in 2022
\cite{mo}, is whether
\[
 \lim_{n\to\infty}\frac{F(n)}{n^{3/2}}
\]
exists.  Here \(A_{ii}=0\), \(A_{ij}=A_{ji}=a_{ij}\in\eps\), and
\[
 Q_A(x)=\sum_{i<j}a_{ij}x_ix_j=\frac12x^{\mathsf T}Ax.
\]

The earlier attempt was rechecked from source rather than assumed.  Its
central conclusions are correct:
\begin{equation}
 \frac1\pi
 \le \liminf_{n\to\infty}\frac{F(n)}{n^{3/2}}
 \le \limsup_{n\to\infty}\frac{F(n)}{n^{3/2}}
 \le\frac12,
 \label{eq:baseline}
\end{equation}
\begin{equation}
 F(n)=\binom n2-2\rho(D_n),
 \qquad
 F(n)\le F(n+1)\le F(n)+n.
 \label{eq:oldfacts}
\end{equation}
The lower bound uses two Gaussian sign roundings and the arcsine law.  The
upper bound uses symmetric Paley conference matrices and principal
submatrices.  The code \(D_n\) is the cut-character set augmented by its
negative.  A claim-by-claim audit, including two attribution errors and a
factor-of-two prose correction, is in \texttt{AUDIT.md}.  Nothing in
\eqref{eq:baseline} or \eqref{eq:oldfacts} implies convergence.

All new statements below distinguish a proved identity from a proposed limit
theorem.  Exact computations guide the search but are not used as asymptotic
proofs.

\section{The relaxed-coupling program collapses}

Let \(m=\binom n2\), let \(c_x=(x_ix_j)_{i<j}\), and for arbitrary real
couplings define
\[
 G_n(a)=\max_{x\in\eps^n}|\langle a,c_x\rangle|.
\]

\begin{theorem}[Cube relaxation]
\label{thm:cube}
The function \(G_n\) is a norm on \(\mathbb R^m\).  Hence
\[
 \min_{a\in[-1,1]^m}G_n(a)=0,
\]
and the unique minimizer is \(a=0\).
\end{theorem}

\begin{proof}
Homogeneity and the triangle inequality follow from the maximum of absolute
linear forms.  If \(X\) is uniform on \(\eps^n\), orthogonality of distinct
degree-two Walsh characters gives
\[
 \E\langle a,c_X\rangle^2=\sum_{i<j}a_{ij}^2.
\]
Thus \(G_n(a)=0\) forces \(a=0\).  Since the origin belongs to the cube, the
last assertion follows.
\end{proof}

In particular, any convex feasible set in the original coupling coordinates
that contains every sign coupling contains their convex hull, the whole cube,
and therefore contains this zero optimizer when the objective remains \(G_n\).
By the audited lower bound, its additive integrality gap is at least
\(n\sqrt{n-1}/\pi\), of the full target order.

The same collapse persists at positive temperature.  Retain duplicated
states \((x,s)\in\eps^n\times\eps\) and put
\[
 L_{n,t}(a)=\frac1t\log\sum_{x,s}
 \exp\bigl(t\langle a,s c_x\rangle\bigr),\qquad t>0.
\]

\begin{theorem}[Finite-temperature cube dual]
\label{thm:cubedual}
The function \(L_{n,t}\) is strictly convex and has the unique cube minimizer
zero, with value
\[
 \min_{a\in[-1,1]^m}L_{n,t}(a)=\frac{(n+1)\log2}{t}.
\]
It has the exact dual representation
\[
 \max_{p\in\Delta(\eps^n\times\eps)}
 \left\{
  \frac{H(p)}t-\left\|\E_p[s c_x]\right\|_1
 \right\}.
\]
\end{theorem}

\begin{proof}
The Hessian of \(L_{n,t}\) is \(t\) times the covariance matrix of \(s c_x\)
under a full-support Gibbs distribution.  If a linear functional of these
vectors were constant, replacing \(s\) by \(-s\) would make the constant zero;
Walsh orthogonality would then make the functional zero.  The Hessian is
therefore positive definite.  Symmetry makes the gradient vanish at zero,
where the value is \(t^{-1}\log 2^{n+1}\).

The entropy variational formula for log-sum-exp gives
\[
 L_{n,t}(a)=\max_p\left\{
 \langle a,\E_p[s c_x]\rangle+\frac{H(p)}t\right\}.
\]
Both domains are compact and convex and the displayed function is convex in
 \(a\) and concave in \(p\), so Sion's theorem applies.  Minimization of a
linear form over the cube is minus the \(\ell^1\)-norm of its coefficient.
\end{proof}

Thus the simplest probability-measure dual is exact but vacuous for the
discrete problem: the uniform state measure cancels every first moment.

There is nevertheless a useful conditional rounding statement.

\begin{proposition}[Variance-sensitive sign rounding]
\label{prop:rounding}
For \(a\in[-1,1]^m\), set \(V(a)=\sum_e(1-a_e^2)\).  There is a sign vector
\(\sigma\in\eps^m\) for which
\[
 G_n(\sigma)\le G_n(a)
 +\sqrt{2V(a)(n+1)\log2}
 +\frac43(n+1)\log2.
\]
\end{proposition}

\begin{proof}
Round the coordinates independently with \(\E\sigma_e=a_e\).  For fixed
\(x\), the centered error
\(\sum_e(\sigma_e-a_e)c_x(e)\) has variance \(V(a)\) and summands bounded by
two.  The two-sided Bernstein bound with parameter \(u\) is at most
\(2e^{-u}\) beyond
\(\sqrt{2V(a)u}+4u/3\).  There are \(2^{n-1}\) distinct cut characters.
Taking \(u=(n+1)\log2\) makes the union failure probability at most \(1/2\).
For a successful outcome, add this uniform error to \(G_n(a)\).
\end{proof}

This gives \(o(n^{3/2})\) loss if \(V(a)=o(n^2)\), but the cube optimizer has
\(V(0)=m\).  The obvious spherical repairs also fail.  On
\(\|a\|_2=\sqrt m\), Parseval gives \(G_n(a)\ge\sqrt m\), attained by a
one-coordinate vector.  Adding \(|a_e|\le1\) to that sphere forces every
\(|a_e|=1\), recovering the original discrete problem exactly.

\section{Temperature scaling and the cross-block wall}

The functional proposed in the question is
\[
 \Phi_{n,\beta}=\min_A\frac1{\beta n^{3/2}}
 \log\sum_x e^{\beta|Q_A(x)|}.
\]
For every fixed \(\beta>0\), the maximum-versus-log-sum-exp inequality gives
\begin{equation}
 \frac{F(n)}{n^{3/2}}
 \le\Phi_{n,\beta}
 \le\frac{F(n)}{n^{3/2}}+\frac{\log2}{\beta\sqrt n}.
 \label{eq:proposedtemp}
\end{equation}
Thus existence of \(\lim_n\Phi_{n,\beta}\) for even one fixed positive
\(\beta\) is equivalent to the original limit.  At this normalization,
fixed temperature is already asymptotically zero temperature.

A genuinely extensive mean-field functional is
\[
 \Psi_{n,\beta}=\min_A\frac1{\beta n}
 \log\sum_{x,s}\exp\left(\frac{\beta sQ_A(x)}{\sqrt n}\right).
\]
It satisfies
\begin{equation}
 \frac{F(n)}{n^{3/2}}
 \le\Psi_{n,\beta}
 \le\frac{F(n)}{n^{3/2}}+\frac{(n+1)\log2}{\beta n}.
 \label{eq:extensivesandwich}
\end{equation}
Consequently, existence of \(\lim_n\Psi_{n,\beta}\) for arbitrarily large
fixed \(\beta\) would settle the problem: the limsup-minus-liminf gap would
be at most \(\log2/\beta\) for every such \(\beta\).

The direct interpolation does not have the needed accuracy.  Define
\[
 h_n(t)=\min_A\log\sum_{x,s}e^{tsQ_A(x)}.
\]

\begin{proposition}[Exact annealed block estimate]
\label{prop:blocktemp}
For all \(n,k\) and \(t>0\),
\[
 h_{n+k}(t)\le h_n(t)+h_k(t)-\log2+nk\log\cosh t.
\]
\end{proposition}

\begin{proof}
For a within-block matrix \(A\), write
\(Z_A^s=\sum_xe^{tsQ_A(x)}\).  Replacing all coefficients of the second
block by their negatives swaps its two one-sided partition sums.  Of the two
possible alignments, one satisfies
\[
 Z_A^+Z_B^+ + Z_A^-Z_B^-
 \le\frac12(Z_A^++Z_A^-)(Z_B^++Z_B^-).
\]
Choose every cross coefficient independently and uniformly.  For fixed
\(x,y,s\), averaging the cross exponential gives \((\cosh t)^{nk}\).
Therefore the expected augmented partition sum is at most the right side
above times this factor, and some deterministic cross signing is no larger.
Take logarithms and infima.
\end{proof}

At \(t=\beta/\sqrt{n+k}\) and \(n/(n+k)\to\alpha\), the normalized cross cost
tends to
\[
 \frac\beta2\alpha(1-\alpha),
\]
which is nonzero for fixed \(\beta\) and grows in the zero-temperature limit.
The direct annealed cross block is therefore a leading term, not an error.

Finally, the absolute value cannot be silently removed.  If every edge from
vertex \(1\) is positive and every other edge negative, fix \(x_1=1\) and put
\(T=\sum_{i>1}x_i\).  Then
\[
 Q_A(x)=\frac n2-\frac{(T-1)^2}{2}.
\]
Hence \(\max_xQ_A(x)=\lfloor n/2\rfloor\), whereas
\(\max_x|Q_A(x)|=\binom n2\).

\section{Covariance relaxations: two exact asymptotic limits}

Let
\[
 \cE_n=\{R\succeq0:\diag R=\mathbf1\}
\]
be the elliptope.  The full arcsine expression is tautologically exact:
\begin{equation}
 M(A)=\max_{R\in\cE_n}
 \left|\frac2\pi\sum_{i<j}a_{ij}\arcsin R_{ij}\right|.
 \label{eq:arcsineexact}
\end{equation}
Indeed, Gaussian sign rounding makes every value on the right at most
\(M(A)\), while \(R=xx^{\mathsf T}\) attains \(Q_A(x)\).  The nonlinearity
therefore preserves the original problem rather than relaxing it.

Linearizing gives a model with a complete thermodynamic limit.  Put
\[
 S(A)=\max_{R\in\cE_n}\left|\frac12\tr(AR)\right|,
 \qquad S_n=\min_A S(A).
\]

\begin{theorem}[Adversarial elliptope limit]
\label{thm:elliptope}
For every \(n\ge2\),
\[
 S_n\ge\frac12n\sqrt{n-1},
\]
and
\[
 \boxed{\displaystyle
 \lim_{n\to\infty}\frac{S_n}{n^{3/2}}=\frac12.}
\]
At each symmetric conference order, equality holds in the finite bound.
\end{theorem}

\begin{proof}
Fix any sign matrix \(A\), set \(q=n-1\), and define
\[
 R^\pm=\frac12\left(I\pm\frac A{\sqrt q}\right)^2.
\]
They are positive semidefinite.  Since every row of \(A\) has squared norm
\(q\), their diagonals equal one.  Moreover
\[
 R^+-R^-=\frac{2A}{\sqrt q}.
\]
The corresponding two linear SDP values differ by
\[
 \frac12\tr\bigl(A(R^+-R^-)\bigr)=n\sqrt q.
\]
At least one has absolute value at least half this amount.

For the reverse asymptotic inequality, if \(R\in\cE_n\), then
\[
 |\tr(AR)|\le\|A\|_{\mathrm{op}}\tr R=n\|A\|_{\mathrm{op}}.
\]
A symmetric conference matrix has operator norm \(\sqrt{n-1}\), so it
attains equality.  For general \(n\), the prime number theorem in arithmetic
progressions supplies a prime \(q\equiv1\pmod4\) with
\(q+1\ge n\) and \(q+1=n(1+o(1))\) \cite{davenport}.  Take a principal \(n\)-by-\(n\)
submatrix of the symmetric Paley conference matrix of order \(N=q+1\).
Compression does not increase its operator norm \(\sqrt{N-1}\), giving
\(S_n\le n\sqrt{N-1}/2=(1/2+o(1))n^{3/2}\).
\end{proof}

Rank-one \(R=xx^{\mathsf T}\) show \(S(A)\ge M(A)\).  The missing transfer is
an asymptotically vanishing integrality gap, which is not supplied by standard
Grothendieck-type constant-factor rounding \cite{alon-naor,nesterov}.  The gap
is already strict at order six: the exact computation below gives \(F(6)=5\),
whereas Theorem~\ref{thm:elliptope} and a conference matrix give
\(S_6=3\sqrt5\).

There is also a sharp barrier for the linear Gaussian-sign method.  Let
\(\Lambda(A)\) be the maximum of
\[
 \sum_{i<j}a_{ij}(R^+_{ij}-R^-_{ij})
\]
over \(R^\pm\in\cE_n\), subject to the edgewise orientations
\(a_{ij}(R^+_{ij}-R^-_{ij})\ge0\).

\begin{theorem}[Oriented covariance barrier]
\label{thm:lambda}
For every sign matrix \(A\),
\[
 M(A)\ge\frac{\Lambda(A)}\pi,
 \qquad
 \Lambda(A)\ge n\sqrt{n-1}.
\]
Moreover
\[
 \boxed{\displaystyle
 \lim_{n\to\infty}
 \frac{\min_A\Lambda(A)}{n^{3/2}}=1.}
\]
\end{theorem}

\begin{proof}
Gaussian sign rounding of \(R^+\) and \(R^-\), followed by the arcsine law,
gives two expected energies.  Since \(\arcsin'(r)\ge1\) and each edge
difference has the prescribed orientation, their difference is at least
\(2\Lambda(A)/\pi\).  Both expectations lie in \([-M(A),M(A)]\), proving the
first inequality.

The square covariances used in Theorem~\ref{thm:elliptope} satisfy the
orientation constraints and have objective \(n\sqrt{n-1}\).  Conversely,
\[
 \Lambda(A)=\frac12\tr\bigl(A(R^+-R^-)\bigr)
 \le n\|A\|_{\mathrm{op}}.
\]
Conference matrices attain equality.  The same prime-number-theorem and Paley
principal-submatrix construction used in Theorem~\ref{thm:elliptope} gives the
matching asymptotic upper bound at every order.
\end{proof}

Thus the optimized two-covariance, edgewise-oriented linear certificate
\(\Lambda\) defined above has adversarial asymptotic lower-bound constant
exactly \(1/\pi\).  Improving it requires information beyond this specific
linearization; the theorem does not exclude more general covariance
certificates.

\section{The augmented cut code and its exact dual}

Encode signs as bits via \((-1)^b\).  Let \(j\in\Ftwo^m\) be the all-one edge
word and
\[
 (\delta u)_{ij}=u_i+u_j.
\]
Then the augmented code in \eqref{eq:oldfacts} is
\[
 D_n=\{\delta u+s j:u\in\Ftwo^n,\ s\in\Ftwo\}.
\]

\begin{theorem}[Linear code and dual]
\label{thm:codedual}
For \(n\ge3\), \(D_n\) is a binary \([m,n]\) linear code and
\[
 D_n^\perp=
 \{H\subseteq E(K_n):
 \deg_H(v)\equiv0\pmod2\ \text{for all }v,
 \ |H|\equiv0\pmod2\}.
\]
For \(n\ge4\), the dual is spanned by 4-cycles, has dimension \(m-n\), and
minimum distance four.
\end{theorem}

\begin{proof}
The cut space has dimension \(n-1\).  Its words have even parity on every
triangle, whereas \(j\) has odd triangle parity, so adjoining \(j\) raises the
dimension by one.  A graph is orthogonal to every cut exactly when every
vertex has even degree, and it is orthogonal to \(j\) exactly when its number
of edges is even.

The cycle space is generated by triangles through one fixed vertex.  An
even-weight cycle-space word is a sum of an even number of these basis
triangles.  The sum of two sharing a nonfixed vertex is a 4-cycle.  Two that
share only the fixed vertex can be paired through an intermediate basis
triangle and expressed as two 4-cycles.  This proves generation.  A nonzero
even-cardinality Eulerian graph has at least four edges, and a 4-cycle exists.
\end{proof}

Counting cuts by the size of one shore gives
\begin{equation}
 W_{D_n}(z)=\frac12\sum_{k=0}^n\binom nk
 \left(z^{k(n-k)}+z^{m-k(n-k)}\right).
 \label{eq:weightenum}
\end{equation}

Let \(a_e=(-1)^{y_e}\) and define the exact augmented partition function
\[
 Z_y(\beta)=\sum_{d\in D_n}
 \exp\bigl(\beta[m-2\wt(y+d)]\bigr).
\]
Its ground state is \(M(A)\).

\begin{theorem}[Signed-Eulerian MacWilliams formula]
\label{thm:macwilliams}
For \(n\ge3\),
\[
 Z_y(\beta)=2^n(\cosh\beta)^mP_y(\tanh\beta),
\]
where
\[
 P_y(t)=\sum_{H\in D_n^\perp}(-1)^{|H\cap y|}t^{|H|}.
\]
For \(0\le t<1\), \(P_y(t)>0\), and uniformly in \(y\),
\[
 M(A)\le\frac1\beta\log Z_y(\beta)
 \le M(A)+\frac{n\log2}\beta.
\]
\end{theorem}

\begin{proof}
The coset MacWilliams identity is
\[
 W_{y+D_n}(z)=2^{n-m}(1+z)^m
 \sum_{H\in D_n^\perp}(-1)^{|H\cap y|}
 \left(\frac{1-z}{1+z}\right)^{|H|}.
\]
Multiply by \(e^{\beta m}\) and substitute \(z=e^{-2\beta}\).  This gives
the formula because
\(e^{\beta}(1+e^{-2\beta})=2\cosh\beta\) and the final ratio is
\(\tanh\beta\).  Positivity follows because the left side is a positive
partition sum.  The last display compares its largest summand with a sum of
\(|D_n|=2^n\) summands.
\end{proof}

For fixed \(\beta>0\), the minimized free energy in
Theorem~\ref{thm:macwilliams}, divided by \(n^{3/2}\), differs from
\(F(n)/n^{3/2}\) by \(O(n^{-1/2})\).  Thus the exact unresolved term is
\[
 \min_y\log P_y(\tanh\beta).
\]
It must cancel the explicit \(m\log\cosh\beta\) term to \(n^{3/2}\)
accuracy.  Averaging \(P_y(t)\) over all \(y\) gives one and yields only
\[
 F(n)\le n\sqrt{(n-1)\log2},
\]
the usual sphere/union-bound constant, weaker than \(1/2\).

Puncturing also stops at the old scale.  For \(n\ge4\), deleting the final star
maps \(D_n\) onto \(D_{n-1}\), and the two lifts of each word have complementary
star coordinates.  Therefore
\[
 \rho(D_{n-1})\le\rho(D_n)
 \le\rho(D_{n-1})+\left\lfloor\frac{n-1}{2}\right\rfloor,
\]
equivalently
\[
 F(n-1)+((n-1)\bmod2)\le F(n)\le F(n-1)+n-1.
\]
The \(n=3\) instance of these inequalities is checked directly; the fiber
there has four rather than two words.
Across a two-block partition, the cross lifts are complementary rank-one
sign matrices.  On the dual side, the
\(\binom n2\binom k2\) alternating cross-block 4-cycles already contribute
macroscopically at mean-field scaling.  Product bounds that take absolute
values discard the cancellation needed here.

\section{High-temperature data do not truncate uniformly}

Let \(X\) be uniform on \(\eps^n\), let \(Y=Q_A(X)\), and use the symmetrized
partition function \(\E\cosh(tY)\), equivalently the augmented state \(sY\).
Direct Walsh counting gives
\begin{equation}
 \E Y^2=m,
 \qquad
 \E Y^4=3m^2-2m+24\sum_{C_4}\prod_{e\in C_4}a_e,
 \label{eq:fourthmoment}
\end{equation}
where the sum is over undirected 4-cycles.  Indeed, a monomial in the fourth
power survives exactly when every vertex has even incidence.  Besides one
edge four times and two edges twice, the only case is a simple 4-cycle, with
multinomial factor \(24\).

Equivalently,
\begin{equation}
 \kappa_4(Y)=3\tr(A^4)-2n(n-1)(3n-4),
 \qquad
 \tr(A^4)=n(n-1)^2+\sum_{i\ne j}(A^2_{ij})^2.
 \label{eq:tracefour}
\end{equation}
The second identity is \(\tr(A^4)=\sum_{ij}(A^2_{ij})^2\); the first follows
by separating the closed walks with four distinct vertices in
\eqref{eq:fourthmoment}.

The expansion
\[
 \log\E\cosh(tY)=\frac{m t^2}{2}+\frac{\kappa_4(Y)t^4}{24}+O(t^6)
\]
shows that, for each fixed \(n\), all sufficiently high-temperature
minimizers of this symmetrized partition function minimize \(\tr(A^4)\).
The symmetrization matters: the one-sided partition can have a
signing-dependent cubic triangle term.  Spectral Cauchy--Schwarz gives
\[
 \tr(A^4)\ge\frac{\tr(A^2)^2}{n}=n(n-1)^2.
\]
For a Paley principal submatrix of order \(n\) inside a conference matrix of
order \(N=n(1+o(1))\),
\[
 \tr(A^4)\le \|A\|_{\mathrm{op}}^2\tr(A^2)
 \le (N-1)n(n-1).
\]
This gives the reverse asymptotic inequality, so
\[
 \lim_n\frac{\min_A\tr(A^4)}{n^3}=1,
 \qquad
 \lim_n\frac{\min_A\kappa_4(Y)}{n^3}=-3.
\]

This fourth-order limit does not determine the next term.  Exhausting all
root-normalized order-seven signings gives
\[
 \min_A\tr(A^4)=342,
\]
while the trace-minimizing signings have three distinct exact sixth moments,
\[
 50{,}781,\qquad 53{,}661,\qquad 59{,}421.
\]
Thus fourth-order data alone are insufficient.  No uniform remainder theorem
connecting this high-temperature expansion to the required ground-state
regime was obtained.

\section{Ordinary graphons and empirical spectra lose the functional}

For a signed matrix \(A\), let \(W_A\) be its step graphon and normalize
\[
 \|W_A\|_\square=\frac1{n^2}
 \max_{S,T\subseteq[n]}
 \left|\sum_{i\in S,j\in T}a_{ij}\right|.
\]

\begin{theorem}[Cut-limit obstruction]
\label{thm:graphon}
If \(M(A_n)=O(n^{3/2})\), then \(W_{A_n}\to0\) in cut norm.  There are two
deterministic sequences satisfying this hypothesis and having the same zero
graphon limit, while their normalized Boolean maxima are separated by at
least \(3/2\).
\end{theorem}

\begin{proof}
For \(x,y\in\eps^n\), split the coordinates according to whether \(x_i=y_i\).
Writing \(z^+,z^-\in\{-1,0,1\}^n\) for the corresponding restrictions of
\(y\), symmetry of \(A\) cancels cross terms and gives
\[
 x^{\mathsf T}Ay=2Q_A(z^+)-2Q_A(z^-).
\]
A multilinear quadratic form attains its maximum absolute value on
\([-1,1]^n\) at a Boolean vertex, so every sign bilinear form is bounded by
\(4M(A)\).  Expressing two indicator vectors as affine combinations of sign
vectors then gives
\[
 \|W_A\|_\square\le\frac{4M(A)}{n^2}.
\]

Take an infinite sequence of Paley conference matrices \(C_N\).  Their
normalized Boolean maxima have limsup at most \(1/2\), and their graphons tend
to zero.  Choose \(k=\lfloor2N^{3/4}\rfloor\) vertices and change every
coefficient inside that set to \(+1\), obtaining \(B_N\).  At most
\(\binom k2=O(N^{3/2})\) edges change, so \(M(B_N)=O(N^{3/2})\) and the
graphon perturbation tends to zero in \(L^1\).  Fix the planted spins to
\(+1\) and average all outside spins independently.  Every cross and outside
term has mean zero, while the planted block contributes \(\binom k2\).
Some outside assignment therefore has at least that energy, and
\[
 \liminf_N\frac{M(B_N)}{N^{3/2}}\ge2.
\]
The separation from the conference limsup is at least \(3/2\).
\end{proof}

The usual dense graph topology therefore collapses every target-scale
near-optimizer to the same point.  This is not repaired by keeping only a
limiting empirical spectrum.  At dense scaling, every empirical law of
\(A/n\) tends to \(\delta_0\), since its average squared eigenvalue is
\((n-1)/n^2\).  At fluctuation scaling there is a stronger obstruction.

\begin{proposition}[Empirical-spectral obstruction]
\label{prop:spectrum}
For every fixed \(d>1/2\), there are conference matrices \(C_N\) and sign
matrices \(B_N\) such that
\[
 W_2(\mu_{B_N/\sqrt N},\mu_{C_N/\sqrt N})\longrightarrow0,
\]
but
\[
 \liminf_N\frac{M(B_N)}{N^{3/2}}\ge d,
 \qquad
 \limsup_N\frac{M(C_N)}{N^{3/2}}\le\frac12.
\]
\end{proposition}

\begin{proof}
Let \(q_N=Q_{C_N}(\mathbf1)\).  The spectral bound gives
\(|q_N|\le N\sqrt{N-1}/2\).  Flip negative coefficients to positive until
the all-one energy is at least \(dN^{3/2}\).  This needs \(O(N^{3/2})\) edges,
and there are \(\Theta(N^2)\) available.  Therefore
\(\|B_N-C_N\|_F^2=O(N^{3/2})\).  Hoffman--Wielandt implies
\[
 W_2^2(\mu_{B_N/\sqrt N},\mu_{C_N/\sqrt N})
 \le\frac{\|B_N-C_N\|_F^2}{N^2}=O(N^{-1/2}).
\]
The energy inequalities follow from the construction and the conference
spectral bound.
\end{proof}

This excludes dense graphons and limiting empirical spectral distributions,
not spectral edges, eigenvectors, or a richer Gaussian-process limit.

For comparison, iid signings are rigorously worse than the conference upper
constant.

\begin{proposition}[Greedy energy of an iid signing]
On one infinite iid Rademacher edge array \(A_n\), almost surely
\[
 \liminf_n\frac{M(A_n)}{n^{3/2}}
 \ge\frac23\sqrt{\frac2\pi}=0.531923\ldots>\frac12.
\]
\end{proposition}

\begin{proof}
Choose spins sequentially by
\[
 x_1=1,\qquad x_j=\operatorname{sgn}\sum_{i<j}a_{ij}x_i.
\]
The resulting energy is
\[
 Q_{A_n}(x)=\sum_{j=2}^n
 \left|\sum_{i<j}a_{ij}x_i\right|=\sum_{k=1}^{n-1}Y_k.
\]
Before column \(j\) is exposed, the earlier spins are measurable with respect
to earlier columns.  Multiplying the fresh iid column signs by those
predictable spins preserves their iid Rademacher law.  Hence the \(Y_k\) are
independent and \(Y_k\) has the law of
\(|\varepsilon_1+\cdots+\varepsilon_k|\).

The central limit theorem and uniform integrability give
\(\E Y_k/\sqrt k\to\sqrt{2/\pi}\).  A Riemann-sum argument therefore yields
\[
 \frac1{n^{3/2}}\sum_{k<n}\E Y_k
 \longrightarrow \sqrt{\frac2\pi}\int_0^1\sqrt t\,dt
 =\frac23\sqrt{\frac2\pi}.
\]
Changing one input sign changes \(Y_k\) by at most two, so its centered moment
generating function is bounded by \(\exp(\lambda^2k/2)\).  Independence gives
an exponential tail \(2e^{-c_\varepsilon n}\) for deviations of
\(\sum_{k<n}Y_k\) by \(\varepsilon n^{3/2}\).  This is summable; Borel--Cantelli
completes the proof.
\end{proof}

This is not a lower bound on \(F(n)\); it shows that typical iid signings are
not asymptotically optimal.

\section{Exact computation through order ten}

Switching makes every edge incident with vertex zero positive.  The remaining
negative coefficients form a graph on \(n-1\) vertices, and vertex
permutations reduce the search to one unlabeled graph in each isomorphism
class.  For a residual negative-edge graph \(G\) and a negative-spin set
\(S\) not containing the root, the exact energy is
\[
 \binom n2-2|E(G)|-2|S|(n-|S|)+4|\delta_G(S)|.
\]
Exhausting nauty's \texttt{geng} streams gives Table~\ref{tab:values}.

\begin{table}[ht]
\centering
\caption{Exact computational values.}
\label{tab:values}
\begin{tabular}{c@{\qquad}rrrrrrrrr}
\toprule
\(n\) & 2&3&4&5&6&7&8&9&10\\
\(F(n)\) & 1&3&4&4&5&9&10&12&13\\
\(F(n)/n^{3/2}\) & .3536&.5774&.5000&.3578&.3402&.4860&.4419&.4444&.4111\\
\bottomrule
\end{tabular}
\end{table}

Completeness is checked against the unlabeled-graph counts and hashes of the
full graph streams.  The reduced formula is checked by direct evaluation;
all labeled signings are independently enumerated through \(n=6\); optional
NetworkX and Z3 routes use different decoders and optimization mechanisms.
Z3 independently gives UNSAT/SAT pairs certifying the displayed values for
\(n=7,8,9\).  At \(n=10\), the lower certificate is the exhaustive nauty
search; the attempted Z3 lower query timed out.

The Paley conference matrix of order ten over \(\mathbb F_9\) has exact
Boolean maximum \(15\), strictly worse than \(F(10)=13\).  The two optimal
switching classes at order ten have different characteristic polynomials.
Thus neither conference structure nor a unique spectrum describes all finite
optimizers.  These finite facts do not imply a limiting constant.

\section{Five proof programs and the sharp remaining lemma}

The detailed derivations and failed variants are preserved in
\texttt{RESEARCH\_LEDGER.md}.  The current programs are logically distinct.

\paragraph{A. Quenched extensive free energy.}
Prove \(\lim_n\Psi_{n,\beta}\) for every fixed \(\beta>0\), or at least for
an unbounded set of \(\beta\)'s.  Equation~\eqref{eq:extensivesandwich} then
forces convergence.  Proposition~\ref{prop:blocktemp} shows exactly why
independent annealed cross blocks fail.  A repair must absorb the cross term
in a joint minimax order parameter, as random-disorder spin-glass
interpolation absorbs it only after disorder expectation
\cite{guerra-toninelli,auffinger-chen}.

\paragraph{B. Near-saturated relaxation.}
Find a nonconvex compact relaxation whose optimizers have
\(\sum_e(1-a_e^2)=o(n^2)\) and an asymptotic composition law.  Then
Proposition~\ref{prop:rounding} transfers its limit to signs.  The cube and
sphere calculations prove that the two obvious feasible sets are structurally
too loose or exactly discrete.

\paragraph{C. Signed Eulerian polynomial.}
Prove a second-order limit for
\(\min_y\log P_y(t)\) after the explicit MacWilliams terms are included.
Puncturing and ordinary products fail because of the common parity phase,
rank-one cross lift, and macroscopic family of mixed 4-cycles.  A quenched
coset-weight theorem, rather than an average weight enumerator, is required.

\paragraph{D. Second-order limit object.}
Construct a rate-sensitive topology for \(A/\sqrt n\) in which the Boolean
maximum is continuous and sign arrays have an asymptotically lossless
composition.  Theorem~\ref{thm:graphon} and
Proposition~\ref{prop:spectrum} rule out ordinary graphons and empirical
spectra.  Spectral edges, eigenvectors, or the whole induced Gaussian process
would have to remain visible.

\paragraph{E. Nonexistence.}
To separate subsequences one needs order-specific lower bounds on \(F(n)\),
not merely constructions with different energies.  Conference, random, and
planted matrices supply upper bounds or deliberately inferior examples.  The
small values fluctuate but do not certify a subsequence mechanism.

No route presently closes.  The single best next lemma is the following.

\begin{center}
\fbox{\begin{minipage}{0.88\textwidth}
\textbf{Quenched minimax free-energy lemma.}
For every fixed \(\beta>0\), the limit
\[
 \lim_{n\to\infty}\min_A\frac1{\beta n}
 \log\sum_{x,s}\exp\left(\frac{\beta sQ_A(x)}{\sqrt n}\right)
\]
exists.
\end{minipage}}
\end{center}

Proving this lemma for an unbounded collection of inverse temperatures settles
the MathOverflow question by \eqref{eq:extensivesandwich}.  In coding
coordinates it asks for cancellation in
\(P_y(\tanh(\beta/\sqrt n))\).  Any proof that bounds the mixed cross-block
contribution independently incurs a leading error; a successful proof must
make that contribution part of the optimized variational structure.

\section*{Reproducibility and final status}

The exact identity checker, exhaustive small-order search, optional Z3
certificate, expected outputs, deterministic seeds, and certificate boundaries
are documented in \texttt{verification/README.md}.  The literature map is in
\texttt{LITERATURE.md}; the original and second proof-search ledgers are both
preserved.

The adversarial elliptope limit, the code dual and MacWilliams formula, the
oriented covariance barrier, and the graph-limit obstructions materially
narrow the search.  They do not prove existence or nonexistence of the
original limit.

\begin{thebibliography}{99}

\bibitem{mo}
P.~Ivanisvili,
\emph{Min-max of a quadratic form of plus-minus ones},
MathOverflow Question 413935, 2022.
\url{https://mathoverflow.net/questions/413935/}

\bibitem{sole-zaslavsky}
P.~Sol\'e and T.~Zaslavsky,
\emph{A coding approach to signed graphs},
SIAM J. Discrete Math. \textbf{7} (1994), 544--553.
\url{https://doi.org/10.1137/S0895480189174374}

\bibitem{guerra-toninelli}
F.~Guerra and F.~L. Toninelli,
\emph{The thermodynamic limit in mean field spin glass models},
Comm. Math. Phys. \textbf{230} (2002), 71--79.
\url{https://doi.org/10.1007/s00220-002-0699-y}

\bibitem{auffinger-chen}
A.~Auffinger and W.-K. Chen,
\emph{Parisi formula for the ground state energy in the mixed p-spin model},
Ann. Probab. \textbf{45} (2017), 4617--4631.
\url{https://doi.org/10.1214/16-AOP1173}

\bibitem{alon-naor}
N.~Alon and A.~Naor,
\emph{Approximating the cut-norm via Grothendieck's inequality},
SIAM J. Comput. \textbf{35} (2006), 787--803.
\url{https://doi.org/10.1137/S0097539704441629}

\bibitem{nesterov}
Y.~Nesterov,
\emph{Semidefinite relaxation and nonconvex quadratic optimization},
Optim. Methods Softw. \textbf{9} (1998), 141--160.
\url{https://doi.org/10.1080/10556789808805690}

\bibitem{lovasz-szegedy}
L.~Lov\'asz and B.~Szegedy,
\emph{Limits of dense graph sequences},
J. Combin. Theory Ser. B \textbf{96} (2006), 933--957.
\url{https://doi.org/10.1016/j.jctb.2006.05.002}

\bibitem{paley}
R.~E. A.~C. Paley,
\emph{On orthogonal matrices},
J. Math. Phys. \textbf{12} (1933), 311--320.
\url{https://doi.org/10.1002/sapm1933121311}

\bibitem{davenport}
H.~Davenport,
\emph{Multiplicative number theory}, third ed.,
Graduate Texts in Mathematics 74, Springer, 2000.
\url{https://doi.org/10.1007/978-1-4757-5927-3}

\end{thebibliography}

\end{document}
